% block: preface-heading page:7 type:section
\section*{{Preface}}\phantomsection\label{sec:preface}
% block: preface-purpose page:7 type:ordinary-paragraph
{In many fields of scientific computing and data science, we frequently face the problem of solving large and sparse linear systems of the form $A\bm x=\bm b$, which is one of the most time-consuming parts of all computations. From this fact, many researchers have devoted themselves to developing efficient numerical algorithms for solving the linear systems, and Krylov subspace methods are nowadays popular numerical algorithms and are known as one of the top ten algorithms of the twentieth century, others including fast Fourier transform and Quick Sort \BookCite{39}. Though the basic theory was established in the twentieth century, Krylov subspace methods have been developed by mathematicians, engineers, physicists, and many others.}
% block: preface-books page:7 type:ordinary-paragraph
{There are many excellent books on Krylov subspace methods, including those by:}
\begin{itemize}
\item Owe Axelsson, 1994 \BookCite{10},
\item Richard Barrett et al., 1994 \BookCite{18},
\item Wolfgang Hackbusch, 1994 \BookCite{92},
\item Are Magnus Bruaset, 1995 \BookCite{29},
\item Rüdiger Weiss, 1996 \BookCite{203},
\item Anne Greenbaum, 1997 \BookCite{81},
\item Yousef Saad, 2003 \BookCite{151},
\item Henk A. van der Vorst, 2003 \BookCite{196},
\item Jörg Liesen {and} Zdeněk Strakoš, 2012 \BookCite{122},
\item Gérard Meurant {and} Jurjen Duintjer Tebbens, 2020 \BookCite{129}.
\end{itemize}
% block: preface-features pages:7,8 type:ordinary-paragraph
{In \BookCite{129}, detailed historical notes of Krylov subspace methods are described, which form a masterpiece (around 700 pages) of Krylov subspace methods for non-Hermitian linear systems. The features of this book are listed as follows:}
\begin{enumerate}[label=(\arabic*)]
\item {Many applications of linear systems from computational science and data science;}
\item {Krylov subspace methods for complex symmetric linear systems such as the COCG method and the COCR method;}
\item {Krylov subspace methods for non-Hermitian linear systems such as the BiCR method, the GPBiCG method, and the (block) IDR($s$) method;}
\item {Krylov subspace methods for shifted linear systems such as the shifted IDR($s$) method;}
\item {Matrix functions as applications of shifted linear systems.}
\end{enumerate}
% block: preface-chapter-two page:8 type:ordinary-paragraph
{Feature (1) corresponds to Chapter~\ref{chap:computational-applications}, and linear systems are derived from various applications: partial differential equations (finite difference discretization methods and the finite element method); computational physics: condensed matter physics (computation of Green's function) and lattice quantum chromodynamics (Wilson fermion matrix); machine learning (least-squares problems); matrix equations (Sylvester-type matrix equations); optimization (Hessian matrix over Euclidean space and Riemannian manifold using tensor computation notations).}
% block: preface-chapter-three page:8 type:ordinary-paragraph
{Features (2) and (3) correspond to Chapter~\ref{chap:krylov-classification}. In this chapter, Krylov subspace methods are classified into three groups: Hermitian linear systems, complex symmetric linear systems, and non-Hermitian linear systems. For Hermitian linear systems, the CG method is derived from the matrix form of the Lanczos process, the CR method is derived from the CG method, and the MINRES method is derived from the Lanczos process. For complex symmetric linear systems, the COCG method, the COCR method, and the QMR\_SYM method are described as extensions of the CG method, the CR method, and the MINRES method, respectively. For non-Hermitian linear systems, the BiCG method, the BiCR method, and the QMR method are described as extensions of the CG method, the CR method, and the MINRES method, respectively. The detailed derivations of the GPBiCG method and the (block) IDR($s$) method, one of the features of this book, are described in this chapter. In addition, some preconditioning techniques are briefly described.}
% block: preface-chapter-four page:8 type:ordinary-paragraph
{Feature (4) corresponds to Chapter~\ref{chap:shifted-linear-systems}. In this chapter, Krylov subspace methods for shifted linear systems are classified into three groups: Hermitian, complex symmetric, and non-Hermitian linear systems. The detailed derivations of these algorithms are described systematically.}
% block: preface-chapter-five page:8 type:ordinary-paragraph
{Feature (5) corresponds to Chapter~\ref{chap:matrix-functions}. If one needs a large matrix function, then Krylov subspace methods and Krylov subspace methods for shifted linear systems are methods of choice since these algorithms can produce any element of the matrix function. The definitions of matrix functions and well-known algorithms for matrix functions are also described.}
% block: preface-numerical-experiments page:8 type:ordinary-paragraph
{An additional feature of this book is that there are no numerical experiments except some typical numerical examples for further understanding the convergence behavior of Krylov subspace methods. The convergence of Krylov subspace methods depends highly on the coefficient matrix, and the best algorithm changes if the coefficient matrix changes. So, if the reader wants to solve linear systems, I recommend the reader to apply several Krylov subspace methods (including the BiCGSTAB method and the GMRES method) to their problem and choose the best one among them.}
% block: preface-conclusion pages:8,9 type:ordinary-paragraph
{This book is suitable for anyone who studied linear algebra and needs to solve large and sparse linear systems. I hope this book is helpful for the reader to understand the principles and properties of Krylov subspace methods and to correctly use Krylov subspace methods to solve their problems.}
\begin{flushright}
{Nagoya, Japan}\qquad Tomohiro Sogabe\\
{September 2022}
\end{flushright}
% block: acknowledgements-heading page:10 type:section
\section*{{Acknowledgements}}\phantomsection\label{sec:acknowledgements}
% block: acknowledgements-colleagues page:10 type:ordinary-paragraph
{I would like to thank Prof. Dr. Lei Du at Dalian University of Technology for kindly sending me the LaTeX code of his block IDR($s$) method. I appreciate Masayuki Nakamura's work. I could raise the priority of writing this book from his timely reminders, together with his kind seasonal greetings. I thank Lokeshwaran Manickavasagam and Poojitha Ravichandran for their help with the production of this book. I am grateful to three anonymous reviewers for their careful reading of the manuscript and highly constructive comments that enhanced the quality of the manuscript. I would like to express my sincere thanks to three editors from the Springer Series in Computational Mathematics (SSCM) editorial board: Prof. Dr. Randolph E. Bank, Prof. Dr. Wolfgang Hackbusch, Prof. Dr. Josef Stoer, and Prof. Dr. Harry Yserentant for providing the author with fruitful comments from broad and distinguished professional viewpoints.}
% block: acknowledgements-family page:10 type:ordinary-paragraph
{Finally, I appreciate Prof. Dr. Shao-Liang Zhang who introduced me to the field of Krylov subspace methods when I was a Ph.D. student at the University of Tokyo about 20 years ago, and I wish to thank my loving wife, Yuan Li, for her day-to-day support and her kindness. She recommended me to write this book on weekends, and she was taking care of two active sons: Kazuma and Takuma, and her old father.}
\clearpage
